%% ehcommands.tex — Custom macros for Event Horizon
%%

% ============================================================
%  Physics — tensor notation
% ============================================================

% Metric tensor
\newcommand{\metric}{g_{\mu\nu}}
\newcommand{\inversemetric}{g^{\mu\nu}}

% Connection coefficients (Christoffel symbols)
\newcommand{\christoffel}[3]{\Gamma^{#1}{}_{#2 #3}}

% Riemann curvature tensor
\newcommand{\riemann}[4]{R^{#1}{}_{#2 #3 #4}}

% Ricci tensor and scalar
\newcommand{\ricci}{R_{\mu\nu}}
\newcommand{\ricciscalar}{R}

% Einstein tensor
\newcommand{\einstein}{G_{\mu\nu}}

% Weyl tensor
\newcommand{\weyl}[4]{C^{#1}{}_{#2 #3 #4}}

% Stress-energy tensor
\newcommand{\stressenergy}{T_{\mu\nu}}

% ============================================================
%  Index notation helpers
% ============================================================

% Partial derivatives
\newcommand{\pu}[1]{\partial^{#1}}           % partial upper index
\newcommand{\pd}[1]{\partial_{#1}}           % partial lower index

% Covariant derivative
\newcommand{\cd}[1]{\nabla_{#1}}

% ============================================================
%  Differential operators
% ============================================================

\newcommand{\laplacian}{\nabla^{2}}
\newcommand{\dalembertian}{\Box}
\newcommand{\lie}[1]{\mathcal{L}_{#1}}      % Lie derivative

% ============================================================
%  ADM decomposition
% ============================================================

\newcommand{\lapse}{\alpha}                  % lapse function
\newcommand{\shift}{\beta^{i}}               % shift vector
\newcommand{\spatialmetric}{\gamma_{ij}}     % spatial 3-metric
\newcommand{\inversespatialmetric}{\gamma^{ij}}
\newcommand{\extrinsic}{K_{ij}}              % extrinsic curvature
\newcommand{\meanextrinsic}{K}               % trace of extrinsic curvature
\newcommand{\spatialchristoffel}[3]{{}^{(3)}\!\Gamma^{#1}{}_{#2 #3}}
\newcommand{\spatialricci}{{}^{(3)}\!R_{ij}}
\newcommand{\spatialricciscalar}{{}^{(3)}\!R}

% ============================================================
%  BSSN / CCZ4 variables
% ============================================================

\newcommand{\conformal}{\chi}                % conformal factor
\newcommand{\conformalmetric}{\tilde{\gamma}_{ij}}
\newcommand{\tracefreecurvature}{\tilde{A}_{ij}}  % trace-free extrinsic curvature
\newcommand{\conformalconnection}{\tilde{\Gamma}^{i}}  % conformal connection functions
\newcommand{\bssnphi}{\phi}                 % BSSN log conformal factor
\newcommand{\bssnW}{W}                      % BSSN W variant

% ============================================================
%  GRMHD
% ============================================================

\newcommand{\lorentz}{W}                     % Lorentz factor
\newcommand{\restmassdensity}{\rho_{0}}
\newcommand{\specificenthalpy}{h}
\newcommand{\magneticfield}[1]{B^{#1}}
\newcommand{\electricfield}[1]{E^{#1}}

% ============================================================
%  Codebase cross-references
% ============================================================

% \coderef{module}{symbol} — margin note linking to code
\newcommand{\coderef}[2]{%
  \marginnote{%
    \textcolor{EHAccretionGold}{\rule{12pt}{1pt}}\\[2pt]%
    {\footnotesize\ttfamily\color{EHMarginGray}%
    #1\linebreak\mbox{}\quad::\,#2}%
  }%
}

% \codemodule{module} — inline monospace module reference
\newcommand{\codemodule}[1]{{\small\ttfamily\color{EHMarginGray}#1}}

% ============================================================
%  Optional sections
% ============================================================

% Crimson star marker for optional/advanced sections
% Safe for use in section/chapter titles (hyperref-compatible)
\newcommand{\starmark}{\texorpdfstring{{\color{EHHorizonCrimson}\(\bigstar\)}}{*}}

% ============================================================
%  Convenience
% ============================================================

% Four-velocity
\newcommand{\fourvel}{u^{\mu}}

% Determinant of metric
\newcommand{\detmetric}{\sqrt{-g}}

% d-dimensional volume element
\newcommand{\vol}[1]{\mathrm{d}^{#1}x}

% Order notation
\newcommand{\order}[1]{\mathcal{O}\!\left(#1\right)}

% Differential
\renewcommand{\d}{\mathrm{d}}

% Absolute value / norm
\newcommand{\abs}[1]{\left\lvert #1 \right\rvert}
\newcommand{\norm}[1]{\left\lVert #1 \right\rVert}

% Inner product
\newcommand{\inner}[2]{\left\langle #1,\, #2 \right\rangle}

\endinput
